\section{Related Work}
\label{sec:related}
\subsection{Coordinated Paths and Predictive Trajectory Planning}
Cooperative planning requires both a compatible passing pattern and motions that vehicles can execute. Preiss et al.~\cite{preiss2017} combined discrete coordination with polynomial refinement and downwash-aware separation. H\"onig et al.~\cite{hoenig2018} used collision-annotated roadmaps to connect multi-agent path finding to continuous quadrotor trajectories. CAPT~\cite{turpin2014} jointly selects assignments and trajectories for interchangeable robots. These approaches exploit discrete coordination to resolve conflicts before continuous refinement. For prescribed assignments and arrival times, the refinement stage must additionally reconcile the selected routes with the available motion budget.

Distributed model predictive control (MPC) resolves interactions through predicted neighbor motions. Luis and Schoellig~\cite{luis2019} introduced on-demand collision constraints and penalized relaxation; Luis et al.~\cite{luis2020} developed online trajectory generation with continuous input parameterization and event-triggered replanning. Buffered Voronoi cells~\cite{zhou2017bvc} instead allocate local free space through relative-position geometry. Communication-efficient predictive formation control~\cite{schiz2025communication} reduces exchanged information using autoencoders. Prediction, local geometric separation, and communication reduction address different computational bottlenecks. Predictive schemes use rolling decision horizons to support online operation, whereas a prescribed common completion time couples the entire remaining motion to its terminal boundary conditions.

\subsection{Safety Filters and Deadlock-Aware Control}
Control barrier functions (CBFs) offer an alternative to repeated full-trajectory planning. Gon\c{c}alves et al.~\cite{goncalves2024} generated drone acceleration fields through a CBF-based quadratic program (QP), followed by thrust and angular-velocity tracking. Their construction provides safety, motion-limit, and QP-feasibility results under its stated conditions. Mestres et al.~\cite{mestre2024} distributed a coupled CBF optimization using auxiliary constraint-mismatch variables, with safety and asymptotic optimality properties for the state-dependent control problem. The optimized quantity in these filters concerns the current control action; a complete trajectory satisfying a prescribed rest-to-rest deadline requires an additional temporal coordination mechanism.

Safe local actions can also encounter coordination deadlocks. Grover et al.~\cite{grover2023} characterized deadlock through the optimality conditions of reactive controllers and developed resolution strategies for two- and three-robot configurations. Jankovic et al.~\cite{jankovic2024} connected gridlock propensity to equilibrium stability and established feasibility for centralized and selected co-optimization policies. Their predictor--corrector strategy illustrates how sharing the responsibility for avoidance improves progress. Chandra et al.~\cite{chandra2025} integrated discrete-time CBFs and deadlock prevention into decentralized receding-horizon optimal control. These developments explicitly address safety and progress, but their navigation formulations do not impose the fixed common rest-to-rest completion task considered here. The relevant distinction is the temporal scope of recovery: local admissible control and eventual progress must be complemented by compatibility with the remaining flight time.

\subsection{Corridors, Motion Limits, and Trajectory Commitment}
Corridors reduce geometric collision restrictions to tractable local constraints. Spatial and relative constructions organize static and mutual avoidance within cooperative optimization~\cite{bu2025}. Park and Kim~\cite{park2021} used relative safe flight corridors to account for uncertain obstacle motion, with slack and reciprocal avoidance to relax their local restrictions. These mechanisms address an important weakness of hard separating corridors: collision-free local regions can still have no intersection with the motions permitted by dynamics and endpoint conditions. A QP solver cannot resolve that geometric incompatibility by optimizing more accurately. Their online formulation updates predictions and avoidance responsibilities; a fixed common arrival time additionally requires recovery of the full rest-to-rest motion within the remaining time allocation. \c{S}enba\c{s}lar et al.~\cite{senbaslar2023} developed RLSS, a networkless planner using linear spatial separations and piecewise-polynomial QPs. RLSS checks derivative limits after optimization and increases segment durations when necessary. Prescribing the completion horizon removes this temporal adjustment and places the recovery burden on motion and geometry.

Online planners also emphasize trajectory commitment and communication. EGO-Swarm~\cite{zhou2021ego} combines local trajectory optimization with perception and implicit topological alternatives. MADER~\cite{tordesillas2022} optimizes trajectories and separating planes with asynchronous check--recheck commitment; Robust MADER~\cite{kondo2024} extends this strategy to communication delays. PRIMER~\cite{kondo2025} learns perception-aware position--yaw planning while retaining deconfliction. Park et al.~\cite{park2025} combine multi-agent waypoints and subgoal optimization, with goal-convergence guarantees in static environments under fully connected communication and deadlock prevention under limited communication. These results establish substantial online coordination capabilities. The present setting focuses on another requirement: returning jointly accepted motions for all assigned UAVs at an unchanged common deadline, including when initial corridor choices conflict with motion constraints.

\subsection{Centralized Cooperative Optimization and Recovery}
Centralized optimization directly couples the motions of interacting vehicles. Adaptive-scaling constrained optimization (ASCO)~\cite{li2021asco} reduces a nonholonomic cooperative planning problem by selecting collision constraints within an adjustable risk range. Failed intermediate solves trigger range contraction, while detected violations expand the active range. The approach therefore includes a recovery mechanism, although its intermediate problems retain nonlinear vehicle dynamics and optimize the terminal time. For the double-integrator task studied here, affine defects and fixed corridor geometry permit convex QPs; box-separation bounds also identify pair restrictions that are unnecessary throughout the current spatial boxes.

Priority-sharing autonomous intersection management~\cite{li2023priority} combines parameterized social-force references, nonlinear cooperative optimization, and batch reservations. Its fixed-duration traverse stage terminates on exit regions, and unsuccessful optimization can lead to a smaller cooperative batch, ultimately a single vehicle. The present recovery task retains the specified team, fixed terminal positions, and rest-to-rest states throughout joint refinement. This boundary-condition commitment concentrates the remaining flexibility in the coordinated motions. It also explains the role of the reference: it supplies a useful passing pattern, while optimization must resolve its remaining incompatibilities.

Sequential convex programming replaces nonlinear separation constraints by local affine restrictions and updates the expansion trajectories~\cite{augugliaro2012}. This yields tractable subproblems, but a chosen linearization can exclude every trajectory compatible with hard motion and endpoint constraints. Enhanced sequential convex programming~\cite{liu2025escp} addresses accumulated linearization errors in predictive formation tracking. In the present model, the integration defects are affine. Softening preserves access to intermediate candidates when current geometric restrictions conflict with exact integration. The fixed-horizon formulation retains this convex structure while endpoint control moments connect recovered candidates to physically checked motion.

\subsection{Fixed-Horizon Feasibility Recovery}
Table~\ref{tab:distinction} compares the decision scope, timing assumptions, and feasibility mechanisms of representative methods. The present task keeps the team, terminal states, and common arrival time fixed. Its central difficulty is compatibility: individually admissible spatial and relative regions can jointly exclude every motion satisfying exact integration. The proposed formulation preserves intermediate candidates through separate kinematic and relative-box penalties, uses them to reconstruct local corridors, and recovers endpoint-consistent controls before complete-interval acceptance. This connects local optimization to the full fixed-horizon task while retaining hard free-space and componentwise motion restrictions.

\begin{table}[htbp]
\begin{adjustwidth}{-\extralength}{0cm}
\caption{Methodological distinctions among representative cooperative planning and safety-control formulations. A prediction horizon specifies a local planning window; a completion horizon specifies the task's arrival time.\label{tab:distinction}}
\small\setlength{\tabcolsep}{4pt}\renewcommand{\arraystretch}{1.15}
\begin{tabularx}{\fulllength}{>{\raggedright\arraybackslash}p{2.65cm}>{\raggedright\arraybackslash}p{3.35cm}>{\raggedright\arraybackslash}p{3.15cm}Y}
\toprule
Method & Coordination and geometry & Time and motion limits & Feasibility or safety mechanism \\
\midrule
ASCO~\cite{li2021asco} & Centralized nonlinear trajectories; selected distance constraints & Optimized terminal time; hard motion bounds & Adapts the active risk range after failed solves or constraint violations; checks the full discretized constraints. \\
\addlinespace[3pt]
Park and Kim~\cite{park2021} & Online polynomial QP; relative safe flight corridors & Receding horizon; derivative bounds in QP & Relative-corridor relaxation and reciprocal avoidance address prediction uncertainty; convex-hull restrictions cover trajectory segments. \\
\addlinespace[3pt]
RLSS~\cite{senbaslar2023} & Decentralized polynomial QP; spatial separating planes & Adjustable durations; derivative limits checked afterward & Geometric separation protects the executed portion; temporal rescaling restores derivative limits when needed. \\
\addlinespace[3pt]
Priority-sharing AIM~\cite{li2023priority} & Centralized nonlinear optimization; social-force initialization & Fixed traverse duration; hard motion bounds & Reservation and batch reduction handle unresolved conflicts; terminal positions lie on exit regions. \\
\addlinespace[3pt]
Gon\c{c}alves et al.~\cite{goncalves2024} & Centralized CBF-QP acceleration field & Instantaneous control; speed and acceleration limits & Feasible safety-filter construction and acceleration tracking under the stated assumptions; no prescribed common arrival time. \\
\addlinespace[3pt]
Mestres et al.~\cite{mestre2024} & Distributed state-dependent CBF optimization & Instantaneous control; waypoint navigation & Constraint-mismatch coordination preserves safety with asymptotic control-optimality properties; no full-horizon endpoint recovery. \\
\addlinespace[3pt]
Jankovic et al.~\cite{jankovic2024} & CBF-QP control; centralized and decentralized policies & Instantaneous acceleration selection & Feasibility analysis and predictor--corrector coordination address collision avoidance and liveness; no common task deadline. \\
\addlinespace[3pt]
Robust MADER~\cite{kondo2024} & Decentralized polynomial planning; separating planes & Receding planning; dynamic limits & Delay-aware trajectory commitment preserves safety despite communication delay. \\
\addlinespace[3pt]
Chandra et al.~\cite{chandra2025} & Decentralized MPC with discrete-time CBFs & Receding horizon; state and input constraints & Liveness-set constraints prevent deadlock in social navigation; safety is imposed through discrete-time barriers. \\
\addlinespace[3pt]
SCP~\cite{augugliaro2012} & Centralized trajectory optimization; local affine separation restrictions & Prescribed trajectory duration; motion constraints & Successive local convex approximations and vehicle-specific feasibility checking. \\
\addlinespace[3pt]
\textbf{Proposed LIOM} & Centralized QP; spatial and relative boxes & Fixed common horizon and endpoints; hard component bounds & Separate penalties preserve intermediate candidates; corridor rebuilding, endpoint control correction, and complete-interval checking recover and accept motion. \\
\bottomrule
\end{tabularx}
\vspace{3pt}
\par\footnotesize CBF: control barrier function; QP: quadratic program; MPC: model predictive control; AIM: autonomous intersection management. Safety statements refer to each formulation's assumptions and execution protocol. The comparison distinguishes mechanisms and task definitions; final acceptance here checks every integrated interval against clearance and motion limits.
\end{adjustwidth}
\end{table}
